\documentclass[a4paper]{article}
\usepackage[pdftex]{graphicx}
\usepackage[utf8]{inputenc}
\usepackage{enumerate}
\usepackage{icomma}
\usepackage{amssymb}
\usepackage{tikz}
\usepackage{href-ul}
\hypersetup{
	colorlinks=true,
	linkcolor=blue,
	urlcolor=blue}
\usepackage{geometry}
\geometry{a4paper, top=15mm, left=15mm, right=15mm, bottom=15mm,
	headsep=10mm, footskip=12mm}

% Start the document
\begin{document}
\begin{tikzpicture}(10,3)
	\draw[ultra thick](2,0) --(10,0) -- (10,1.3) --(2,1.3) -- (2,0);
	\draw[fill=black](2,0)-- (0,.6) -- (2,.6) -- (2,0);
	\node at (6,.6) {\large Lösungsblatt von \href{https://www.okuyakl.de}{www.okuyakl.de}};
\end{tikzpicture}
\vspace{0.5 cm}

\renewcommand{\arraystretch}{2.8}
\begin{tabular}{|p{40 pt}|p{50 pt}|p{80 pt}|p{100 pt}|p{100 pt}|p{50 pt}|}
\hline
Funktion & Öffnung $a$ & Scheitel $(x|y)$ & Funktionsterm $f(x)=$&  Normalform $f(x)=$ & Anzahl Nullstellen\\
\hline
$f_1(x)$  & $-3$ & $(-8|5)$ & $-3(x+8)^2+5$ & $ -3x^2-48x-187$ & 2\\
\hline
$f_2(x)$  & $-2$ & $(-4|0)$ & $-2(x+4)^2$ & $-2x^2-16x-32 $ & 1 \\
\hline
$f_3(x)$  & $0,5$ & $(-2|2)$ & $0,5(x+2)^2+2 $ & $0,5x^2+2x+4$ & 0 \\
\hline
$f_4(x)$  & $0,1$ & $(0|-2,5)$ & $0,1x^2-2,5$ & $0,1x^2-2,5$& 2 \\ 
\hline
$f_5(x)$  & $-0,25$ & $(2|5)$ & $-0,25(x-2)^2+5$ & $-0,25x^2+x+4$ & 2 \\
\hline
$f_6(x)$  & $4$ & $(4|-6)$ & $4(x-4)^2-6$ & $4x^2-32x + 58$ & 2 \\ 
\hline
$f_7(x)$  & $0,01$ & $(6|-4)$ & $0,01(x-6)^2-4$  & $0,01x^2-0,12x-3,64$ & 2\\ 
\hline

\hline
\end{tabular}

\begin{center}
	\includegraphics[width=7 cm]{../../viecher/endcomic.pdf}
	
	Hier geht es zurück zum \href{https://www.okuyakl.de/math/m9aparL016/aa016.pdf}{Aufgabenblatt}
\end{center}

\end{document}